\begin{papercard}{lamonica2021}
\cf{Stages} \stagetag{D} \stagetag{C} \stagetag{W} \stagetag{H}. A three-phase study: co-design workshops on how users want privacy information presented, a privacy policy \emph{risk assessment tool} applied to 34 apps/e-tools, and consensus privacy guidelines for clinicians.

\cf{Aims} (1) Consumer preferences for accessing privacy information; (2) develop and apply a privacy policy risk assessment tool against recommended standards; (3) propose guidelines for health professionals \pc{lamonica2021}{1}{This study aims to present consumer preferences for accessing privacy information}.

\cf{Phase 1: co-design} Ten participatory design workshops (July--September 2018) in nine headspace youth mental-health centres, with young people aged $\geq$12, supportive others and professionals \pc{lamonica2021}{2}{Our research team conducted a series of 10 participatory design}. The workshops ran in three stages: discovery, evaluation of platform images, and prototyping with sketchbooks \pc{lamonica2021}{3}{stages: discovery, evaluation, and prototype}. Knowledge translation: two young researchers with no prior knowledge of the platform reviewed all notes and artefacts independently and identified key concepts \pc{lamonica2021}{3}{reviewed by an independent knowledge translation}. Ethics approval: University of Sydney HREC 2018/130.

\cf{Phase 2: instrument} Built from published digital-health privacy/security criteria. The APA app evaluation model was the primary reference, supplemented by the Australian National Safety and Quality Digital Mental Health Standards \pc{lamonica2021}{3}{Although the American Psychiatric}. The article describes 23 items in four sections: policy (6), personal health information (4), data security and storage (9), other aspects (4). Validity is claimed as face/construct validity because the items come from standards, but \emph{no validity analysis} was done \pc{lamonica2021}{3}{we were confident that the measure has}. The appendix numbers 25 questions with skip logic and mostly nominal or graded answer sets \pp{lamonica2021tool}. Examples: ``Is it clear how and when the organisation will delete personal health information?'' (not clear / somewhat clear / very clear) \pc{lamonica2021tool}{2}{Is it clear how and when the organisation will delete personal health}; recipients typed as relevant vs.\ irrelevant third parties, academic partners, government, health-related people \pc{lamonica2021tool}{2}{Relevant third parties (e.g. for analytic purposes, customer support}; retention categories (unspecified, until no longer needed, $<$1 year, 1--3 years, $>$3 years, at user discretion) \pc{lamonica2021tool}{3}{UNTIL NO LONGER NEEDED}; breach response explained; explicit warning about risks of services outside the organisation's control \pc{lamonica2021tool}{4}{Does the privacy policy explicitly warn users about privacy risks}.

\cf{Material capture} Sample: the 34 apps/e-tools in the youth configuration of the InnoWell platform, a convenience sample of a curated platform rather than store-based. Policies \emph{and related material} such as terms were taken from websites and store links with a ``nonrestrictive process'' that followed all hyperlinks from the store page or within the policy \pc{lamonica2021}{4}{nonrestrictive process was used, allowing all hyperlinks from}. Assessment was done primarily by one team member with a research intern \pc{lamonica2021}{4}{developed primarily}, and no agreement statistic is reported.

\cf{Readability} Flesch Reading Ease and Flesch--Kincaid grade, computed on only the \emph{first two paragraphs} (at least 100 words) of each policy \pc{lamonica2021}{4}{calculated using the first two paragraphs of the privacy policies}. The appendix instructs using Word's readability statistics on the first paragraph \pc{lamonica2021tool}{1}{Please copy and paste the first paragraph from the privacy policy}. This sampling choice limits validity.

\cf{Phase 3: guidelines} Consultations in 2014--2016 and 2017--2018 with stakeholder organisations, including two former privacy commissioners. Eight core foundation principles: legal compliance, anonymous/pseudonymous use where practicable, individual control, transparency, interaction with individuals, encryption, de-identification, cross-border risk \pc{lamonica2021}{4}{core foundation principles included responsibility for legal}. The final checklist was agreed by research-team consensus \pc{lamonica2021}{4}{resulting in agreement by consensus}. It is given as Textbox 1, phrased as questions a clinician asks \pc{lamonica2021}{9}{Privacy guidelines—health professional and service provider considerations}.

\cf{Key findings} Workshop participants wanted privacy information \emph{before account creation}, plus granular control over who sees what \pc{lamonica2021}{4}{participants wanted privacy information to be presented}. All 34 had a policy; 91\,\% were rated college-graduate level and 9\,\% professional level \pc{lamonica2021}{6}{All but three policies fell into the college graduate grade level}. Deletion timing was unclear in 44\,\% \pc{lamonica2021}{6}{Although more than half of the privacy policies}. Only 1 app allowed changing sharing preferences, and only by email \pc{lamonica2021}{6}{of apps allowed data sharing preferences to be changed}. 76\,\% did not explain breach response, and only 24\,\% warned of risks outside the provider's control \pc{lamonica2021}{8}{Less than one-quarter of privacy policies}.

\cf{Limitations stated} The tool and guidelines were developed by the research team without the broader stakeholder group, legal or industry experts, or people with lived experience \pc{lamonica2021}{11}{we acknowledge that the privacy policy risk assessment}.

\cf{Use for this thesis} The most directly reusable items for \emph{protection framing and limits} (RQ3): ``cannot guarantee'' style risk warnings, breach response, retention categories, and relevant vs.\ irrelevant third parties. The graded clarity answers (not / somewhat / very clear) show how to code vagueness. Weak points to avoid: a convenience sample of a curated platform, single-assessor scoring, readability from two paragraphs only, and no validation.
\end{papercard}


\begin{papercard}{lamonica2021tool}
\cf{What it is} The 4-page multimedia appendix of LaMonica et al.\ (2021): the form of the ``Privacy Policy Risk Assessment Tool'' \pc{lamonica2021tool}{1}{Privacy Policy Risk Assessment Tool}. The article reports 23 items; the form numbers 25 (the target audience and reading-level questions are extra), and the card follows the form. Four blocks: privacy policy (Q1--Q7), personal health information (Q8--Q11), data security and storage (Q12--Q21), other aspects (Q22--Q25). Skip logic is written into the form (``If Yes Q3; If No Q8'').

\cf{Policy block} Q1 asks for the target audience as free text \pc{lamonica2021tool}{1}{Who is the target audience (age range or population) for the app or e-tool}. Q3 measures reading level: the rater pastes the policy's first paragraph into Word and notes the Flesch reading ease and Flesch--Kincaid grade \pc{lamonica2021tool}{1}{note the Flesch reading ease and Flesch-kincaid grade level}. Q4 asks whether the app claims to meet the Australian Privacy Act 1988, HIPAA or an international equivalent \pc{lamonica2021tool}{1}{Does the app or e-tool claim to meet the standards of the Privacy Act 1988 (Australia), HIPAA}; Q5--Q7 ask whether the policy introduces its purpose, the organisation and ``adequate'' information for all users, with answers no / some / yes.

\cf{Personal health information and storage} Q8 asks whether personal health information is collected \pc{lamonica2021tool}{2}{Does the app or e-tool collect personal health information}; Q9 whether it is shared; Q10 and Q11 are clarity ratings (not clear / somewhat clear / yes, very clear) for security steps \pc{lamonica2021tool}{2}{Is it clear if the organisation has taken reasonable steps to ensure the security of personal health}. Q12 asks whether use and purpose are declared \pc{lamonica2021tool}{2}{Does the privacy policy declare data use and purpose}. Q13 types recipients as A relevant third parties, B ``irrelevant'' third parties such as social media \pc{lamonica2021tool}{2}{Irrelevant third parties (e.g. Social media)}, C academic partners, D government, E health-related people and F other. Q14 de-identification \pc{lamonica2021tool}{3}{Is shared data de-identified}; Q15 sharing preferences \pc{lamonica2021tool}{3}{Can the user change their preferences regarding data sharing}; Q16 deletion \pc{lamonica2021tool}{3}{Can the user delete data from the app or e-tool}; Q17--Q18 on-device versus server storage \pc{lamonica2021tool}{3}{Is user data stored on the device}; Q19 retention with ordinal categories (unspecified, until no longer needed, under 1 year, 1--3 years, over 3 years, at user discretion) \pc{lamonica2021tool}{3}{For how long is data stored}; Q20 server type and Q21 country \pc{lamonica2021tool}{3}{What type of server is used to store user data}.

\cf{Other aspects} Q22 third-party vendors such as analytics \pc{lamonica2021tool}{4}{Does the app use third party vendors}; Q23 breach response \pc{lamonica2021tool}{4}{Is the manner in which the organisation will respond to a data breach adequately explained}; Q24 complaint and enquiry routes, scored no / some (internal \emph{or} external) / yes (both) \pc{lamonica2021tool}{4}{Does the privacy policy inform users as to how they can make enquires, provide feedback}.

\cf{Limits (card)} Several items need judgement without anchors (``adequate information'' Q7, ``irrelevant'' third parties Q13B, ``somewhat clear'' Q10--Q11); Q24 combines two conditions in one question; the legal and storage references are Australian; readability is computed from a short extract with Word's English-language formula, so it cannot be transferred to German policies; the form states no rater procedure or agreement check.

\cf{Use for this thesis} A compact list of policy-coding topics with typed answer sets: recipient typology, retention categories, storage location, deletion, breach plan, complaint routes. For the codebook, add observable anchors for the clarity levels (for example what counts as ``somewhat clear'' for deletion), replace the Australian references with GDPR and DiGA terms, and keep the skip logic so that coders do not rate items that do not apply.
\end{papercard}
